\documentclass[9pt,twoside]{pnas-new}
\articletype{SOCIAL SCIENCES}
\templatetype{pnasmathematics}
% Remove DRAFT watermark
\setboolean{displaywatermark}{false}
\begin{document}
\title{Coal dependence does not independently predict renewable energy support:
       Evidence from survey data, production records, and news framing analysis}
\author[a,1]{Alexandra S. Popescu}
\affil[a]{Program in Quantitative Social Science, Dartmouth College,
          Hanover, NH 03755}
\leadauthor{Popescu}
\significancestatement{Public resistance to renewable energy is often attributed
to community dependence on fossil fuels. Using a cross-sectional survey
(\textit{N}\,=\,2,400), sixty years of state-level coal production records, and
zero-shot natural language processing applied to 122 news articles, this study
finds that coal dependence does not independently predict renewable energy support
once individual ideology and climate concern are controlled. However, coal
dependence amplifies conservative resistance, creating a compounded disadvantage
in coal communities. These findings suggest that structural economic interventions,
not communication campaigns, are the more promising path to building durable
public support for the energy transition.}
\authorcontributions{A.S.P. merged all data, conducted all statistical analyses and natural language processing, and wrote the manuscript.}
\authordeclaration{The author declares no competing interests.}
\correspondingauthor{\textsuperscript{1}To whom correspondence should be
addressed. E-mail: alexandra.s.popescu.28@dartmouth.edu}
\keywords{renewable energy support $|$ coal dependence $|$ political ideology
          $|$ energy transition $|$ zero-shot NLP}
\begin{abstract}
The energy transition requires broad public support, yet communities with deep
ties to coal often exhibit the greatest skepticism. Whether that skepticism is
rooted in place, specifically in the economic history of fossil fuel production,
or in the individual political attitudes that cluster in those places is an open
question with direct implications for policy design. This study addresses that
question by combining individual-level survey data (\textit{N}\,=\,2,400),
state-level coal production records from the U.S. Energy Information Administration
(1960--2023), power plant proximity measures, and a zero-shot natural language
processing analysis of 122 news articles. Five models are estimated: ordinary
least squares (OLS) regression, an OLS model with a coal share $\times$ ideology
interaction, a multilevel model (MLM) with state random intercepts, a CatBoost
gradient-boosted tree model with Shapley Additive Explanations (SHAP), and a
zero-shot news framing classifier using BART-large-MNLI. All five consistently
show that state-level coal dependence does not independently predict renewable
energy support once individual ideology and climate concern are controlled
(intraclass correlation $\approx$\,0.00). However, a significant
coal\,$\times$\,ideology interaction ($\hat{\beta} = -0.035$, $p = .049$) reveals
that conservative residents of coal-dependent states face a compounded disadvantage
beyond what ideology alone predicts. News framing analysis yields no significant
association between coal share and climate-focused coverage
($\chi^2 = 5.18$, $p = .159$). These findings suggest that resistance to
renewable energy in fossil fuel communities is driven primarily by individual
political identity, amplified by economic context. Structural economic
interventions, rather than information campaigns, are the more promising lever
for building the coalitions that a managed energy transition requires.
\end{abstract}
\dates{This manuscript was compiled on \today}
\doi{\url{www.pnas.org/cgi/doi/10.1073/pnas.XXXXXXXXXX}}
\maketitle
\thispagestyle{firststyle}
\ifthenelse{\boolean{shortarticle}}{\ifthenelse{\boolean{singlecolumn}}%
  {\abscontentformatted}{\abscontent}}{}
\Firstpage
\section*{Introduction and Related Work}
The United States faces a structural paradox in its energy transition: renewable
energy commands majority public support nationally, yet the communities with the
deepest ties to fossil fuel production often remain its most vocal opponents.
Policy analysts and communication scholars have long debated whether this
resistance reflects the economic and structural features of coal-dependent places,
or whether it is more parsimoniously explained by the individual political
identities that happen to concentrate in those communities \cite{Stokes2020,
Mildenberger2019, Carley2020}. The answer matters considerably for policy design.
If place shapes attitudes independently of individual ideology, then targeted
regional interventions, whether economic or communicative, may shift opinion. If
place is merely a proxy for the people who live there, then such investments may
yield little return.

Prior research establishes strong individual-level predictors of renewable energy
support. Political ideology consistently emerges as the strongest negative
predictor, while environmental concern and educational attainment are positive
predictors \cite{Ansolabehere2014, Hess2014, Krupnick2010}. Less settled is
whether community-level fossil fuel dependence exerts an independent effect once
these individual characteristics are controlled. Some scholarship argues that
carbon-intensive communities develop distinctive political cultures that suppress
clean energy support through economic self-interest, identity formation, and local
media framing \cite{Mildenberger2019, Jacobs2016}. Others contend that place
adds little explanatory power once the compositional characteristics of residents
are measured \cite{Carley2020, Ansolabehere2014}.

This study addresses the debate directly through a multi-method empirical design
that combines survey data, administrative production records, and computational
text analysis. I test three related propositions: first, whether state-level coal
dependence predicts renewable energy support after controlling for individual
ideology and climate concern; second, whether coal dependence moderates the
relationship between ideology and support, specifically whether conservative
residents of coal states face a compounded disadvantage; and third, whether news
media in coal-dependent states frame energy coverage differently in ways that might
constitute a proximate mechanism linking place to attitudes.

\section*{Data and Methods}
\subsection*{Survey Data}
Individual-level data come from a cross-sectional survey of U.S. adults
(\textit{N}\,=\,2,400) conducted in 2024. The dependent variable is renewable
energy support (Q137), measured on a five-point Likert scale
(1\,=\,strongly oppose to 5\,=\,strongly support), with a mean of 3.96
(SD\,=\,0.94). Predictors include political ideology (1--7, liberal to
conservative), climate concern (composite), educational attainment, income,
age, gender, and race/ethnicity drawn from a supplementary respondent file.
After listwise deletion the analytic sample comprises
\textit{N}\,=\,2,389 respondents nested in 48 states.

\subsection*{State-Level Coal Dependence}
State coal production data come from the U.S. Energy Information Administration's
State Energy Data System (SEDS), 1960--2023 \cite{EIA2023}. Coal share is defined
as each state's average share of total energy production attributable to coal over
1990--2023, capturing the economic legacy of fossil fuel dependence rather than
any single year's snapshot.

\subsection*{Power Plant Proximity}
Plant-level locations from EIA Form EIA-860 are used to compute state-level
counts of active coal-fired and utility-scale solar and wind facilities,
included as additional covariates to capture exposure to incumbent and emerging
energy infrastructure.

\subsection*{Regression Models}
All continuous predictors are standardized (mean\,=\,0, SD\,=\,1) before
estimation. The baseline OLS model regresses renewable energy support on
individual and state-level predictors. The interaction model adds a
coal\,$\times$\,ideology term. The mixed-effects model includes a random
intercept by state estimated via restricted maximum likelihood; the ICC
quantifies variance attributable to between-state differences. The swing-state
model replicates the baseline specification on the 12 most contested states
in the 2024 presidential election ($n = 687$).

\subsection*{CatBoost and SHAP}
A CatBoost gradient-boosted tree ensemble \cite{Prokhorenkova2018} was estimated
with hyperparameters tuned via five-fold cross-validation (depth, learning rate,
number of iterations). Variable importance was assessed using SHapley Additive
exPlanations (SHAP) \cite{Lundberg2017}, which decompose each prediction into
additive feature contributions.

\subsection*{Zero-Shot News Framing}
One hundred twenty-two news articles were collected via NewsAPI
(January 2023 to March 2024) using search terms combining ``renewable energy,''
``coal,'' and U.S. state names. Each article was classified using zero-shot
natural language inference with \texttt{facebook/bart-large-mnli} \cite{Lewis2020}
via the HuggingFace \texttt{transformers} library, with candidate labels:
\textit{climate}, \textit{economic benefit}, \textit{energy security},
and \textit{other}. The highest-probability label was assigned to each article.
A Pearson chi-square test assessed whether frame distributions differed between
high- and low-coal states (split at the sample median). A 50-state OLS regression
tested the continuous association between coal share and percentage climate framing.

\Endparasplit

\section*{Results}

\subsection*{Renewable Energy Support Is High Across the Political Spectrum}

Before turning to the regression models, Figure~\ref{fig:statebar} presents mean
renewable energy support by state, colored by the 2024 presidential election
winner. Three patterns are immediately visible. First, support is high and
relatively uniform: 46 of 48 states in the analytic sample record mean RE support
above the scale midpoint of 3.0, and 44 exceed 3.5. Second, the partisan gap,
while present, is modest. Democratic-won states cluster near the top of the
distribution, but a large number of Republican-won states record means at or above
the sample mean of 3.96. Third, the clearest outliers are Wyoming (mean\,$\approx$\,2.4)
and Oklahoma (mean\,$\approx$\,3.5), both high-coal, Republican-won states, which
hints at a role for fossil fuel context that the descriptive pattern alone cannot
disentangle from ideology.

Critically, several 2024 swing states, including Michigan, Wisconsin, Pennsylvania,
Georgia, Arizona, North Carolina, and Nevada, all of which were won by the
Republican candidate, record mean RE support close to or above the national mean.
This suggests that even in the most electorally contested, politically mixed
environments, renewable energy enjoys broad baseline support, and that the
variation within Republican-won states is as informative as the variation between
partisan groups. This descriptive picture motivates the regression models that
follow: if place explained attitudes, we would expect Republican-won coal states
to cluster distinctly below Democratic-won states; the fact that they do not
points toward individual-level drivers.

\begin{figure}[h]
\centering
\includegraphics[width=\linewidth]{re_support_ranked_by_winner}
\caption{Mean renewable energy support (1--5) by state, ranked from highest to
lowest and colored by 2024 U.S. presidential election winner
(blue\,=\,Democratic, red\,=\,Republican). The dashed vertical line marks the
sample mean (3.96). Swing states that were carried by the Republican candidate
in 2024, including Michigan, Wisconsin, Pennsylvania, Georgia, Arizona, North
Carolina, and Nevada, are highlighted in red yet record mean RE support near or
above the national average, consistent with the regression finding that place
does not independently predict attitudes (\textit{N}\,=\,2,389).}
\label{fig:statebar}
\end{figure}

\subsection*{Coal Dependence Does Not Independently Predict Renewable Energy Support}

Table~\ref{tab:models} presents standardized regression coefficients across four
model specifications. In the baseline OLS model, climate concern is the strongest
positive predictor of renewable energy support
($\hat{\beta} = +0.275$, $p < .001$), followed by educational attainment
($\hat{\beta} = +0.068$, $p < .001$) and income
($\hat{\beta} = +0.045$, $p < .01$). Political ideology is the strongest
negative predictor ($\hat{\beta} = -0.129$, $p < .001$), with more conservative
respondents reporting substantially lower support. Together, individual-level
predictors explain approximately 21\% of variance in renewable energy support
($R^2 = .211$).

State-level coal share does not reach statistical significance in any additive
specification ($\hat{\beta} = -0.018$, $p = .34$ in the baseline OLS;
$\hat{\beta} = -0.022$, $p = .59$ in the swing-state subsample). This null
finding is robust across the multilevel model, the swing-state subgroup, and
the machine learning specification described below. Once individual ideology
and climate concern are accounted for, a state's coal production history
provides no independent predictive power over residents' attitudes toward
renewable energy. The near-uniformly high RE support observed in Figure~\ref{fig:statebar},
even across Republican-won swing states, is consistent with this result.

The multilevel model confirms that this null finding is not an artifact of
model misspecification. The intraclass correlation coefficient approaches zero
(ICC\,$\approx$\,0.00), indicating that virtually no variance in renewable energy
support is attributable to between-state differences after individual predictors
are included. The random state intercept adds negligible explanatory value,
reinforcing the conclusion that renewable energy attitudes are driven by who
people are rather than where they live.

\begin{table*}[t!]
\centering
\caption{Standardized regression coefficients predicting renewable energy support
(1--5) across four model specifications. Standard errors in parentheses.}
\label{tab:models}
\begin{tabular*}{\textwidth}{@{\extracolsep{\fill}}lcccc@{}}
 & OLS Baseline & OLS + Interaction & MLM & Swing States \\
\midrule
\multicolumn{5}{l}{\textit{Individual-level predictors}} \\
Climate concern
  & $0.275^{***}$ & $0.274^{***}$ & $0.274^{***}$ & $0.261^{***}$ \\
  & $(0.021)$     & $(0.021)$     & $(0.021)$     & $(0.038)$     \\
Ideology (conservative)
  & $-0.129^{***}$ & $-0.131^{***}$ & $-0.130^{***}$ & $-0.144^{***}$ \\
  & $(0.019)$      & $(0.019)$      & $(0.019)$      & $(0.034)$      \\
Education
  & $0.068^{***}$ & $0.068^{***}$ & $0.068^{***}$ & $0.071^{**}$ \\
  & $(0.018)$     & $(0.018)$     & $(0.018)$     & $(0.033)$    \\
Income
  & $0.045^{**}$ & $0.045^{**}$ & $0.045^{**}$ & $0.039$ \\
  & $(0.018)$    & $(0.018)$    & $(0.018)$    & $(0.033)$ \\
\multicolumn{5}{l}{\textit{State-level predictor}} \\
Coal share
  & $-0.018$ & $-0.019$ & $-0.017$ & $-0.022$ \\
  & $(0.019)$ & $(0.019)$ & $(0.022)$ & $(0.041)$ \\
\multicolumn{5}{l}{\textit{Interaction}} \\
Coal share $\times$ Ideology
  & & $-0.035^{*}$ & & \\
  & & $(0.018)$    & & \\
\midrule
$R^2$ / ICC & $.211$ & $.213$ & $\approx .000$ & $.198$ \\
$N$         & $2{,}389$ & $2{,}389$ & $2{,}389$ & $687$ \\
\bottomrule
\end{tabular*}
\addtabletext{All continuous predictors are standardized (mean\,=\,0, SD\,=\,1).
$^{*}p < .05$;\; $^{**}p < .01$;\; $^{***}p < .001$.
MLM $R^2$ is reported as ICC (proportion of variance due to state-level clustering).
Swing state subsample includes respondents from the 12 most contested states in
the 2024 U.S. presidential election.}
\end{table*}

\subsection*{Coal Dependence Amplifies Ideological Resistance}

The OLS model with an interaction term reveals a significant
coal\,$\times$\,ideology effect ($\hat{\beta} = -0.035$, $p = .049$). Among
liberal respondents (ideology at $-$1\,SD), coal share has a near-zero marginal
effect on renewable energy support. Among conservative respondents (ideology at
$+$1\,SD), residing in a high-coal state is associated with meaningfully lower
support, a pattern that is invisible in additive models that omit the interaction
term. Figure~\ref{fig:marginal} illustrates this divergence: as coal share
increases, predicted support declines most sharply among self-identified
conservatives, producing a ``double disadvantage'' for those who are both
ideologically conservative and economically embedded in fossil fuel production.
Wyoming, annotated in Figure~\ref{fig:marginal} as a high-coal conservative
state, exemplifies this pattern: it is the only state in Figure~\ref{fig:statebar}
with mean RE support below 2.5, a level that cannot be explained by ideology
alone given that other heavily Republican states without significant coal
dependence record means close to the national average.

This interaction is the study's primary substantive finding. It suggests that
the political resistance to clean energy visible in states like West Virginia or
Wyoming is not simply a product of conservative ideology in isolation, but of an
amplification process by which economic context reinforces ideological
predispositions. This pattern has important implications for both theory and
policy, as discussed below.

\begin{figure}[h]
\centering
\includegraphics[width=\linewidth]{marginal_effects_interaction}
\caption{Marginal effect of coal share on predicted renewable energy support by
political ideology. Predicted values are derived from the OLS model with a
coal\,$\times$\,ideology interaction term; all other covariates are held at
their means. Shaded bands represent 95\% confidence intervals. Liberal and
conservative are defined as ideology at $-$1\,SD and $+$1\,SD from the mean,
respectively (\textit{N}\,=\,2,389).}
\label{fig:marginal}
\end{figure}

\subsection*{SHAP Attribution Confirms Primacy of Individual Predictors}

The gradient-boosted tree model (CatBoost) confirms the dominance of
individual-level predictors in a non-parametric framework that makes no linearity
assumptions. SHAP analysis identifies climate concern as the single largest
contributor to predicted renewable energy support, followed by political ideology,
educational attainment, and income. Coal share ranks among the lowest-importance
features in the SHAP attribution plot, consistent with the null OLS finding and
providing convergent validity from a machine learning perspective. The non-linear
model uncovers no meaningful threshold effects or higher-order interactions beyond
the coal\,$\times$\,ideology term already identified in the parametric models.

\subsection*{News Framing Does Not Vary with Coal Dependence}

Zero-shot classification of 122 news articles yields the following frame
distribution: climate-focused (38\%), economic benefit (31\%), energy security
(19\%), and other (12\%). A chi-square test comparing high- and low-coal states
(split at the sample median) finds no significant difference in frame
distribution ($\chi^2 = 5.18$, df\,=\,3, $p = .159$). As shown in
Figure~\ref{fig:framing}, low-coal states receive somewhat more climate-focused
coverage (47\% vs. 28\%) and high-coal states somewhat more economic framing
(36\% vs. 26\%), but these differences do not reach conventional significance
thresholds. Extending the analysis to all 50 states using a continuous OLS
regression of percentage climate framing on coal share likewise yields a null
result ($\hat{\beta} = -0.08$, $p = .42$, $R^2 = .01$). Media coverage of energy
issues does not appear to vary systematically with state coal dependence, ruling
out differential news framing as a proximate mechanism linking place to attitudes.

\begin{figure}[h]
\centering
\includegraphics[width=\linewidth]{framing_bar}
\caption{Distribution of news frames by state coal dependence group. States are
split at the sample median coal share into high-coal and low-coal groups.
Bars show the percentage of articles classified into each frame by the
BART-large-MNLI zero-shot classifier. Although low-coal states show a higher
share of climate framing and high-coal states a higher share of economic framing,
the difference in frame distributions is not statistically significant
($\chi^2 = 5.18$, df\,=\,3, $p = .159$; \textit{N}\,=\,122 articles).}
\label{fig:framing}
\end{figure}

\section*{Discussion}

\subsection*{Place Does Not Drive Renewable Energy Attitudes}

The central finding of this study is a null result of substantive importance:
state-level coal dependence does not independently predict renewable energy
support once individual ideology and climate concern are controlled. Figure~\ref{fig:statebar}
makes this visible at a glance: Republican-won swing states such as Michigan,
Wisconsin, and Pennsylvania, states that carry significant electoral weight and
mixed economic histories, record mean RE support near or above the national
average of 3.96. The intraclass correlation from the multilevel model approaches
zero, and coal share is consistently non-significant across OLS, swing-state, and
machine learning specifications. This suggests that the apparent link between
fossil fuel communities and clean energy skepticism is largely a compositional
effect, a reflection of who lives in coal states rather than the states themselves.

This finding challenges policy frameworks premised on the assumption that place
shapes attitudes through economic self-interest or regional political culture.
Targeted information campaigns about local renewable energy benefits, a common
intervention design, are unlikely to move attitudes that are primarily a product
of individual-level ideological identity. The appropriate unit of intervention
is the person, not the place.

\subsection*{The Interaction Finding and Its Policy Implications}

The one exception to this pattern is both statistically significant and
theoretically meaningful. The coal\,$\times$\,ideology interaction
($\hat{\beta} = -0.035$, $p = .049$) reveals that living in a coal-dependent
state amplifies ideological resistance to renewable energy, and this effect
holds only among conservatives. This ``double disadvantage,'' illustrated in
Figure~\ref{fig:marginal}, is invisible in additive models and may explain why
prior research has reached conflicting conclusions about the role of place in
shaping energy attitudes. Wyoming's position as the lowest-ranked state in
Figure~\ref{fig:statebar}, despite not being dramatically more Republican than
Kansas or Oklahoma, is consistent with this interaction: high coal share
compounds the ideological effect in ways that additive models cannot capture.

The implication is that energy transition policy must do more than change minds.
In communities where conservative ideology and coal dependence interact to
suppress renewable energy support, durable change is unlikely to emerge from
communication strategies alone. Structural economic interventions are the more
promising lever for building the necessary coalitions, including job guarantee
programs, severance and retraining funds, and community benefit agreements tied
to renewable siting \cite{Carley2020, Stokes2020}. These interventions operate
on the material conditions that amplify ideological resistance, rather than on
ideology itself.

\subsection*{Null Framing Result as a Finding}

The absence of a significant relationship between coal share and news framing,
shown in Figure~\ref{fig:framing}, rules out one plausible mechanism: that
coal-dependent states receive disproportionately economic or energy-security
coverage that dampens climate concern and renewable energy support. Although the
direction of the framing differences is consistent with theoretical expectations,
the magnitude is too small to reach significance at the sample sizes available.
This shifts the explanatory burden back toward political socialization,
economic self-interest, and identity processes operating largely outside direct
media influence, suggesting that framing interventions targeted at coal-state
media markets would be unlikely to succeed on their own.

\subsection*{Limitations and Future Directions}

Several limitations bound these conclusions. The survey is cross-sectional,
precluding causal inference. Coal share is measured at the state level, which
may mask important county-level variation in fossil fuel dependence and economic
exposure. The news corpus (\textit{N}\,=\,122) is modest and limited to
English-language national outlets, potentially missing hyper-local media most
salient to community attitude formation. Zero-shot classification, while
reproducible and scalable, may not capture nuanced framing distinctions that
fine-tuned models or human coders would identify. Future work should use panel
survey data, county-level economic exposure measures, and larger and more
geographically targeted media corpora to adjudicate among the attitude formation
mechanisms this study cannot resolve.

\section*{Agentic Analysis}

Artificial intelligence tools played an integral role at every stage of this
project, and QSS~45's explicit encouragement to engage with these technologies
critically and purposefully shaped how I approached them. The first substantive
use came during the literature review phase, where I used NotebookLM as a
reading and synthesis partner. Uploading twelve sources and querying the tool
across them simultaneously allowed me to map the conceptual structure of the
field far more efficiently than reading sequentially would have permitted. The
mind map NotebookLM generated made the dominant clusters visible at a glance:
the Energy Justice framework, the structural vulnerability literature, and the
framing debate. What the exercise also made clear, however, was the importance
of verification. When I cross-checked three specific claims the tool generated
against the original sources, none was fully accurate: two were partially
accurate, presenting only part of the evidence, and one was a hallucination
that attributed a characterisation of Appalachian jobs as ``high-wage'' to
Graff and Carley (2018) when the source made no such claim. This taught me
that source-grounded AI reduces hallucination but does not eliminate it, and
that the researcher remains responsible for every claim that carries their name.

For the quantitative and computational components of the project, I used
Claude as an agentic coding assistant, which required a different set of skills
than simply querying a chatbot. I learned that the quality of the output
is directly proportional to the precision of the prompt. Vague instructions
such as ``run a regression'' produced generic code; specifying the dependent
variable, the standardisation procedure, the exact interaction term, and the
desired output format produced working, interpretable results. I also learned
to prompt for explanation alongside execution: asking the model not only to
produce code but to explain what each step does and why forced me to engage
with the methodological logic rather than treat the output as a black box.
This was essential for interpreting results correctly, particularly for the
multilevel model's ICC and the SHAP feature attributions, where understanding
what the number means conceptually is as important as obtaining it. The
zero-shot NLP pipeline required especially careful prompt design: specifying
the candidate labels, the model, and the aggregation logic step by step
produced reproducible classifications, whereas underspecified prompts returned
inconsistent frame assignments. Across all of these tasks, the course's
emphasis on understanding method before delegating execution proved essential.
AI can accelerate analysis, but only a researcher who understands what OLS
assumes, why an interaction term changes the interpretation of main effects,
and what a null ICC implies can evaluate whether the output is correct.
The most important skill QSS~45 gave me is not knowing how to use these tools,
but knowing when and whether to trust them.

\dataavail{Survey data are available upon request. EIA SEDS coal production data
are publicly available at \url{https://www.eia.gov/state/seds/}. EIA Form EIA-860
plant data are available at \url{https://www.eia.gov/electricity/data/eia860/}.
Analysis code and processed datasets are archived at
\url{https://github.com/alexandrapopescu/coal-re-support}.}

\acknow{The author thanks Professor Herbert Chang at Dartmouth College for
guidance throughout this project. All analyses were conducted in Python~3.10
using \texttt{statsmodels}, \texttt{pymer4}, \texttt{catboost}, \texttt{shap},
and \texttt{transformers}.}

\showacknow

\begin{thebibliography}{99}

\bibitem{Stokes2020}
Stokes LC (2020)
\newblock \textit{Short Circuiting Policy: Interest Groups and the Battle Over
Clean Energy and Climate Policy in the American States}
(Oxford Univ Press, New York).

\bibitem{Mildenberger2019}
Mildenberger M (2019)
\newblock \textit{Carbon Captured: How Business and Labor Control Climate Politics}
(MIT Press, Cambridge, MA).

\bibitem{Carley2020}
Carley S, Konisky DM (2020)
\newblock The justice and equity implications of the clean energy transition.
\newblock \textit{Nat Energy} 5(8):569--577.

\bibitem{Ansolabehere2014}
Ansolabehere S, Konisky DM (2014)
\newblock \textit{Cheap and Clean: How Americans Think About Energy in the Age of
Global Warming}
(MIT Press, Cambridge, MA).

\bibitem{Hess2014}
Hess DJ (2014)
\newblock Sustainability transitions: A political coalition perspective.
\newblock \textit{Res Policy} 43(2):278--283.

\bibitem{Krupnick2010}
Krupnick A, et al. (2010)
\newblock Public views on a 21st-century electricity policy.
\newblock \textit{Public Opin Q} 74(3):449--473.

\bibitem{Jacobs2016}
Jacobs M (2016)
\newblock The politics of economic growth.
\newblock \textit{J Policy Hist} 28(1):113--136.

\bibitem{EIA2023}
U.S. Energy Information Administration (2023)
\newblock State Energy Data System (SEDS): 1960--2023.
\newblock Available at \url{https://www.eia.gov/state/seds/}.
Accessed January 2024.

\bibitem{Prokhorenkova2018}
Prokhorenkova L, Gusev G, Vorobev A, Dorogush AV, Gulin A (2018)
\newblock CatBoost: Unbiased boosting with categorical features.
\newblock \textit{Adv Neural Inf Process Syst} 31:6638--6648.

\bibitem{Lundberg2017}
Lundberg SM, Lee SI (2017)
\newblock A unified approach to interpreting model predictions.
\newblock \textit{Adv Neural Inf Process Syst} 30:4765--4774.

\bibitem{Lewis2020}
Lewis M, et al. (2020)
\newblock BART: Denoising sequence-to-sequence pre-training for natural language
generation, translation, and comprehension.
\newblock \textit{Proc 58th Annu Meet Assoc Comput Linguist}, pp 7871--7880.

\end{thebibliography}

\end{document}
